\documentclass[12pt,a4paper]{article}
% \usepackage[utf8]{inputenc} % fontspecと併用時は不要
\usepackage[japanese]{babel}
\usepackage{amsmath}
\usepackage{amsfonts}
\usepackage{amssymb}
\usepackage{geometry}
\usepackage{graphicx}
\usepackage{enumitem}
\usepackage{fancyhdr}
\usepackage{verbatim}

% 日本語改行設定
\usepackage{xeCJK}
\usepackage{indentfirst}
\setlength{\parindent}{1em}

% 日本語フォント設定
\usepackage{fontspec}
\setmainfont{Hiragino Mincho Pro}
\setsansfont{Hiragino Sans}
\setmonofont{Menlo}
\setCJKmainfont{Hiragino Mincho Pro}
\setCJKsansfont{Hiragino Sans}
\setCJKmonofont{Menlo}
% ページ設定
\geometry{margin=1.5cm}
\pagestyle{fancy}
\fancyhf{}
\fancyhead[L]{プログラミング応用 第三回}
\fancyhead[R]{演習問題解説}
\fancyfoot[R]{\ifnum\value{page}=1\small（裏面に続く）\fi}
\fancyfoot[C]{\thepage}
\renewcommand{\headrulewidth}{0.4pt}
\renewcommand{\footrulewidth}{0.4pt}

% 問題番号のカウンター
\newcounter{problem}
\setcounter{problem}{0}

% シンプルな問題環境
\newenvironment{problem}[1][]{
    \refstepcounter{problem}
    \vspace{1em}
    \noindent
    \textbf{問\arabic{problem}} #1
    %\vspace{0.5em}
    %\par
}{
    \vspace{1em}
}

% 解答環境
\newenvironment{solution}{
    \vspace{0.5em}
    \noindent
    \textbf{【解答】}
    \vspace{0.3em}
}{
    \vspace{1em}
}

\begin{document}

% ヘッダー情報
\begin{center}
    {\Large \textbf{プログラミング応用 第三回}}\\[0.5em]
    {\large \textbf{演習問題解説}}\\[1em]
\end{center}

\begin{problem}
次の問いに答えよ.
\begin{enumerate}
\item $k$個のボールそれぞれを, $k$個の箱の中から一様ランダムに選んで入れる. このとき, どの箱にもちょうど一つずつボールが入る確率を求めよ.
\item 1で求めた確率を $p_k$ とする. $\lim_{k\to\infty} \frac{\log p_k}{k}$ の極限値を求めよ (ただしここでは自然対数を考える).
ただし, 解答では次のStirlingの近似を証明なしに用いてよい:
\begin{align*}
    \sqrt{2\pi} n^{n+1/2} e^{-n} \le n! \le e n^{n+1/2} e^{-n}
\end{align*}
\item $k$個のボールを$10k^2$個の箱の中から一つ一様ランダムに選んで入れる. このとき, どの箱にもボールが高々一つしか入らない確率が少なくとも$0.9$であることを示せ.
\end{enumerate}
\end{problem}


\begin{solution}
\begin{enumerate}
\item $p_k = \frac{k!}{k^k}$.
\item $\lim_{k\to\infty} \frac{\log p_k}{k} = -1$
\item $1\le i<j\le k$に対し
\[
\text{$i$番目と$j$番目のボールが同じ箱に入る}
\]
という事象を$\mathcal{E}_{i,j}$とする. このとき, $\Pr[\mathcal{E}_{i,j}] = \frac{1}{10k^2}$である. 「全ての箱に高々一つのボールが入る」の余事象は「ある$1\le i<j\le k$が存在して, $\mathcal{E}_{i,j}$が成り立つ」である. この確率は高々
\begin{align*}
  \Pr\left[\bigvee_{1\le i<j\le k} \mathcal{E}_{i,j}\right] \le \sum_{1\le i<j\le k} \Pr[\mathcal{E}_{i,j}] = \frac{k(k-1)}{2} \cdot \frac{1}{10k^2} \le \frac{1}{10}
\end{align*}
である. 従って, 所望の確率は少なくとも$0.9$である.
\end{enumerate}
\end{solution}

\newpage
\begin{problem}
    ベクトル $v\in\mathbb{R}^n \setminus\{0\}$ を固定する.
    ベクトル $r \in\{0,1\}^n$ を, 各$r_i$を独立に確率$1/2$で$0$または$1$にする.
    このとき
    \begin{align*}
        \Pr\left[\sum_{i=1}^n v_i r_i \ne 0\right] \ge 1/2
    \end{align*}
    が成り立つことを示せ.
\end{problem}

\begin{solution}
\begin{itemize}
  \item $v$が非ゼロベクトルなので, ある$i\in\{1,\dots,n\}$が存在して, $v_i\ne 0$である. 特に$v$の成分の順番を入れ替えても確率は不変なので, $v_1\ne 0$としても一般性を失わない.
  \item 考えている確率を書き換えると
  \begin{align*}
    \Pr\left[\sum_{i=1}^n v_i r_i \ne 0\right] &= \Pr\left[v_1r_1 + \sum_{i=2}^n v_i r_i \ne 0\right] \\
    &= \Pr\left[v_1r_1 \ne - \sum_{i=2}^n v_i r_i\right] \\
    &= \Pr\left[r_1 \ne -\frac{1}{v_1} \sum_{i=2}^n v_i r_i \right]
  \end{align*}
  \item $r_2,\dots,r_n$の値が何であっても, $r_1$が$-\frac{1}{v_1} \sum_{i=2}^n v_i r_i$と一致する確率は高々$1/2$である. 従って, 所望の確率は少なくとも$1/2$である.
\end{itemize}
\end{solution}

\newpage

\begin{problem}
確率変数$X,Y$をそれぞれ$\{0,\dots,5\}$上の一様分布に従う確率変数とする ($X$と$Y$は独立).
このとき, 確率変数の独立性の定義に基づいて以下の問いに答えよ.
\begin{enumerate}
\item $X$ と $(X+Y)\bmod 6$ は独立であるかどうか, 理由も含めて答えよ.
\item $X$ と $(X\cdot Y)\bmod 6$ は独立であるかどうか, 理由も含めて答えよ.
\end{enumerate}
\end{problem}

\begin{solution}
採点は答えだけ確認すればよい (証明はおわなくてよい)
\begin{enumerate}
\item 正しい.
\item 誤り.
\end{enumerate}
\end{solution}

\newpage

\begin{problem}
ベクトル $y\in\mathbb{R}^D$ を固定する.
行列$B\in\mathbb{R}^{d\times D}$を, 各成分が独立に標準正規分布 $N(0,1)$ に従うランダム行列とし,
$A=\frac{1}{\sqrt{d}}B$ とする.
このとき, $\mathbb{E}[\|Ay\|^2] = \|y\|^2$ が成り立つことを示せ (講義資料に書いてある内容は証明なしに用いて良い).

\end{problem}

\begin{solution}
\begin{itemize}
  \item まず, $Ay$の代わりに$By$のノルムの期待値を考える. $(By)_i=\sum_{j=1}^D B_{ij} y_j$であるから,
  \begin{align*}
    \mathbb{E}[(By)_i^2] &= \mathbb{E}\left[\left(\sum_{j=1}^D B_{ij} y_j\right)^2\right] \\
    &= \mathbb{E}\left[\sum_{1\le j,j'\le D} B_{ij}B_{ij'} y_j y_{j'}\right] \\
    &= \sum_{j=1}^D y_j^2 \mathbb{E}[B_{ij}^2] + \sum_{j\ne j'} y_j y_{j'} \mathbb{E}[B_{ij}]\mathbb{E}[B_{ij'}] \\
    &= \sum_{j=1}^D y_j^2 \\
    &= \|y\|^2.
  \end{align*}
  \item これを$i=1,\dots,d$に対して足し合わせると, $\mathbb{E}[\|By\|^2] = d \|y\|^2$ が得られる.
  \item 両辺を$d$で割ると所望の等式が得られる.
\end{itemize}
\end{solution}
\end{document}
